\begin{table}[htbp]
  \centering
  \caption{Exact power of a complete case analysis when the two dropout probabilities differ and each group is inflated by its own dropout probability. Design parameters: $\pi_{1} = 0.45$, $\pi_{0} = 0.30$, $\rho_{1} = \rho_{0} = 0$, $r = 1$, $\alpha = 0.025$ (one-sided) and target power $= 0.8$. The target power is attained throughout, so equal dropout probabilities are not necessary for a complete case design to reach its target.}
  \label{tab:unequal-dropout}
  \begin{threeparttable}
  \begin{tabular}{ccrrrrr}
    \toprule
    $\omega_{1}$ & $\omega_{0}$ & $n_{\mathrm{complete}}$ & $n_{1}$ & $n_{0}$ & Power & Type I error \\
    \midrule
    0.10 & 0.10 & 163 & 182 & 182 & 0.8048 & 0.0251 \\
    0.05 & 0.20 & 163 & 172 & 204 & 0.8035 & 0.0251 \\
    0.20 & 0.05 & 163 & 204 & 172 & 0.8035 & 0.0251 \\
    0.02 & 0.35 & 163 & 167 & 251 & 0.8036 & 0.0252 \\
    0.35 & 0.02 & 163 & 251 & 167 & 0.8036 & 0.0251 \\
    0.40 & 0.05 & 163 & 272 & 172 & 0.8033 & 0.0252 \\
    \bottomrule
  \end{tabular}
  \begin{tablenotes}[flushleft]\footnotesize
    \item $n_{\mathrm{complete}}$ is the number of completers per group the unadjusted formula requires, and $n_{j} = \lceil n_{\mathrm{complete}} / (1 - \omega_{j}) \rceil$ is the number randomized to group $j$. Power and type I error rate are exact, computed by complete enumeration over the numbers of completers and responders rather than by simulation.
  \end{tablenotes}
  \end{threeparttable}
\end{table}
